%% FullCourse_10Lecs -- Lecture 6, Topic 6.4: Deep Equilibrium Nets + OLG Aggregate Uncertainty
%% Source: ./archive_pre_split/Lec06_HA_Models.tex (sections 8-9)
%% Pass B additions: Validation frames (Euler residuals + economic plausibility) ported from MiniCourse M2_T4 sec 5.
%% FullCourse-only: OLG with Aggregate Uncertainty section (not present in MiniCourse).

%% Shared preamble for the FullCourse_10Lecs topic decks.
%%
%% Each topic .tex \input's this file, then defines its own \title, \date
%% override (if any), \graphicspath, and \begin{document}.
%%
%% Reconstructed verbatim from the inlined copy preserved in
%% Lec10_Case_Studies/Lec10_T8_Case6_DeepRamsey.tex (lines 23-190), which is
%% explicitly annotated there as "from ../shared_preamble.tex, copied verbatim".
%% Base preamble matches MiniCourse_8hr/Slides/shared_preamble.tex; this version
%% additionally provides \sectiondivider, the conceptbox/cautionbox/examplebox/
%% takeawaybox tcolorboxes, and \compacttables used across the split decks.

\documentclass[10pt,english,aspectratio=169]{beamer}

\usepackage{ae,aecompl}
\usepackage[T1]{fontenc}
\usepackage{textcomp}
\usepackage[utf8]{inputenc}
\usepackage{babel}
\usepackage{datetime2}

\setcounter{secnumdepth}{3}
\setcounter{tocdepth}{3}

\usepackage{amsmath, amssymb, amsthm, graphicx}
\usepackage{booktabs, multirow}
\usepackage{tikz}
\usepackage{pgfplots}
\pgfplotsset{compat=1.16}
\usepackage[most]{tcolorbox}
\tcbuselibrary{skins, breakable, theorems}
\usepackage{listings}
\usepackage{xcolor}

\lstset{
  basicstyle=\ttfamily\small,
  keywordstyle=\color{blue!80!black}\bfseries,
  commentstyle=\color{green!50!black}\itshape,
  stringstyle=\color{orange!80!black},
  backgroundcolor=\color{gray!8},
  frame=single,
  rulecolor=\color{gray!40},
  breaklines=true,
  showstringspaces=false,
  numbers=none,
}

\lstdefinestyle{pythonstyle}{
    language=Python,
    basicstyle=\ttfamily\footnotesize,
    keywordstyle=\color{blue!80!black}\bfseries,
    commentstyle=\color{green!50!black}\itshape,
    stringstyle=\color{orange!80!black},
    frame=single,
    breaklines=true,
    showstringspaces=false,
    numbers=left,
    numbersep=5pt,
    numberstyle=\tiny\color{gray},
    backgroundcolor=\color{gray!8},
    rulecolor=\color{gray!40}
}

\ifx\hypersetup\undefined
  \AtBeginDocument{%
    \hypersetup{bookmarks=false,breaklinks=true,pdfborder={0 0 0},colorlinks=true,
      linkcolor=DarkText,urlcolor=RedTitle}%
  }
\else
  \hypersetup{bookmarks=false,breaklinks=true,pdfborder={0 0 0},colorlinks=true,
    linkcolor=DarkText,urlcolor=RedTitle}%
\fi

\makeatletter
\newcommand\makebeamertitle{%
  \begin{frame}[plain]
    \vspace*{1.0cm}
    \centering
    {\usebeamerfont{title}\usebeamercolor[fg]{title}\inserttitle\par}
    \vspace{1.0cm}
    {\usebeamerfont{author}\usebeamercolor[fg]{author}\insertauthor\par}
    \vspace{0.2cm}
    {\usebeamerfont{date}\usebeamercolor[fg]{date}\insertdate\par}
  \end{frame}}%
\AtBeginDocument{%
  \let\origtableofcontents=\tableofcontents
  \def\tableofcontents{\@ifnextchar[{\origtableofcontents}{\gobbletableofcontents}}
  \def\gobbletableofcontents#1{\origtableofcontents}%
}

\usetheme{metropolis}
\usefonttheme{professionalfonts}

\definecolor{LightGray}{RGB}{230,230,230}
\definecolor{RedTitle}{RGB}{0,0,200}
\definecolor{VeryLightGray}{RGB}{245,245,245}
\definecolor{DarkText}{RGB}{0,0,0}

\setbeamercolor{background canvas}{bg=white}
\setbeamercolor{title}{fg=RedTitle}
\setbeamercolor{author}{fg=DarkText}
\setbeamercolor{date}{fg=DarkText}
\setbeamercolor{frametitle}{bg=VeryLightGray,fg=DarkText}
\setbeamercolor{normal text}{fg=DarkText}
\setbeamerfont{title}{size=\large}

\setbeamersize{text margin left=5mm, text margin right=5mm}
\addtobeamertemplate{frametitle}{}{\vspace{-0.5em}}
\newcommand{\reducevspace}{\vspace{-0.5cm}}

% Shared section divider and box styles for split topic decks.
\newcommand{\sectiondivider}[2]{%
\begin{frame}[plain,noframenumbering]
\begin{center}
\vfill
{\Large\textbf{#1}\par}
\vspace{0.35cm}
{\normalsize #2\par}
\vfill
\end{center}
\end{frame}
}

\newtcolorbox{conceptbox}[1][]{
  colback=blue!3,
  colframe=RedTitle!55,
  boxrule=0.35mm,
  arc=1mm,
  left=2mm,
  right=2mm,
  top=1mm,
  bottom=1mm,
  #1
}
\newtcolorbox{cautionbox}[1][]{
  colback=red!3,
  colframe=red!50!black,
  boxrule=0.35mm,
  arc=1mm,
  left=2mm,
  right=2mm,
  top=1mm,
  bottom=1mm,
  #1
}
\newtcolorbox{examplebox}[1][]{
  colback=green!3,
  colframe=green!45!black,
  boxrule=0.35mm,
  arc=1mm,
  left=2mm,
  right=2mm,
  top=1mm,
  bottom=1mm,
  #1
}
\newtcolorbox{takeawaybox}[1][]{
  colback=VeryLightGray,
  colframe=RedTitle!55,
  boxrule=0.35mm,
  arc=1mm,
  left=2mm,
  right=2mm,
  top=1mm,
  bottom=1mm,
  #1
}
\newcommand{\compacttables}{\renewcommand{\arraystretch}{1.15}}

\setbeamertemplate{navigation symbols}{}
\setbeamertemplate{footline}{%
  \hfill%
  \usebeamercolor[fg]{page number in head/foot}%
  \usebeamerfont{page number in head/foot}%
  \insertframenumber\,/\,\inserttotalframenumber\kern1em\vskip2pt%
}

\makeatother

\author{Zhigang Feng}
% "date with time": compile timestamp so each build is identifiable.
\date{\today\ \DTMcurrenttime}


\graphicspath{{./pic/}{../../MiniCourse_8hr/Slides/Module2_DL_RL_Macro/pic/}{../}}

\title[AI for Econ Research]{AI for Economic Research: Dynamic Models, Language, and Agents\\
Lecture 6.4: DEQN and OLG with Aggregate Uncertainty}

\begin{document}

\makebeamertitle

%=============================================================================
\begin{frame}{Topic 6.4 Agenda}
\begin{enumerate}
  \item \textbf{Deep Equilibrium Nets (DEQN)}
  \medskip
  \item \textbf{OLG with Aggregate Uncertainty}
\end{enumerate}
\medskip
\textit{Arc: T1 Aiyagari + computation $\to$ T2a Why ML / Euler $\to$ T2b Global DNN $\to$ T3 KS / DeepHAM $\to$ \textbf{T4 DEQN / OLG} $\to$ T5 MFG $\to$ Lab06}
\end{frame}

%=============================================================================
\section{Deep Equilibrium Nets (DEQN)}

\begin{frame}{Deep Equilibrium Nets (Azinovic, Gaegauf \& Scheidegger 2022, \emph{IER})}
\begin{itemize}
\item DEQN is a \textbf{general} method: a \textbf{single} policy network, trained on equilibrium conditions, no value function. The original paper introduced it on \textbf{OLG / Bewley} economies.
\medskip
\item \textit{In this lecture} we \textbf{apply} DEQN to the KS economy (the Lab06 setup). For KS the single network represents the policy:
\[
\sigma_\theta(a, e, Z, K) \in (0,1)
\;\Rightarrow\;
a' = \underline a + \sigma_\theta\cdot(\text{cash on hand} - \underline a).
\]
\medskip
\item The training objective is the \textbf{Euler equation residual}, not a Bellman residual.
\medskip
\item No value network, no fictitious play; just supervised learning on the Euler condition.
\medskip
\item Key innovation: a \emph{cloud} of parallel economies provides a rich training distribution.
\end{itemize}

\begin{tcolorbox}[colback=blue!3, colframe=blue!40, title={Why DEQN is the capstone of the deep-HA toolbox}]
\scriptsize DEQN combines the Euler-loss idea (Topic 6.2a) with the policy-network architecture (Topic 6.2b) --- without ever computing a value function. Azinovic et al.\ (2022) demonstrate it on OLG/Bewley models; \emph{this lecture} applies the same recipe to the full KS heterogeneous-agent benchmark.
\end{tcolorbox}
\end{frame}

\begin{frame}{DEQN: policy network architecture}
\begin{itemize}
\item Single multilayer perceptron $\sigma_\theta : \mathbb{R}^4 \to (0,1)$.
\medskip
\item Inputs: $(a, e, Z, K)$ where $K = \int a\, d\mu$ is aggregate capital.
\medskip
\item Architecture applied in this lecture (Lab06):
  \begin{itemize}
  \item 3 hidden layers, 64 units each.
  \item Mish activation function.
  \item Sigmoid output for the savings rate.
  \end{itemize}
\medskip
\item Sigmoid output enforces the budget constraint at the network level: the implied $a'$ always lies within $[\underline a, \text{cash on hand}]$.
\medskip
\item Contrast with DeepHAM: DeepHAM uses both $V_{\text{NN}}$ and $\mathcal C_{\text{NN}}$; DEQN uses only the policy network.
\end{itemize}
\end{frame}

\begin{frame}{DEQN: Euler equation as the loss function}
The household optimality condition is the Euler equation:
\[
u'(c_t) \;=\; \beta\,\mathbb{E}_t\!\left[(1+r_{t+1})\,u'(c_{t+1})\right].
\]
\medskip
The \textbf{Euler residual}:
\[
\varepsilon \;=\; 1 \,-\, \frac{\beta\,\mathbb{E}\!\left[(1+r')\,u'(c')\right]}{u'(c)}.
\]
\medskip
The DEQN loss aggregates squared residuals across a sample of agents:
\[
\mathcal{L}(\theta) \;=\; \frac{1}{M}\sum_{m=1}^{M} \varepsilon_m^2.
\]
\medskip
\textbf{Why Euler equations?} They are \emph{necessary} for optimality; sufficiency requires in addition a transversality / no-Ponzi condition (plus concavity). DEQN works because the bounded state space and simulated paths effectively impose transversality, so driving $\varepsilon\to 0$ trains the optimal policy without ever computing $V$.
\end{frame}

\begin{frame}{DEQN: cloud simulation}
\textbf{Problem.} To evaluate the Euler residual we need $(K, Z)$, but $K$ depends on the cross-sectional distribution $\mu$, which itself evolves endogenously under the trained policy.
\medskip

\textbf{Solution.} Simulate $N$ \emph{independent copies} of the entire economy in parallel:
\medskip
\begin{itemize}
\item Epoch 0: all $N$ economies start at the steady-state distribution $\mu_{ss}$ from a sequence-Jacobian (SSJ) calibration.
\item As training progresses, each economy evolves under its own TFP path $\{Z_n\}$, so the set $\{K_n\}$ spreads out across a realistic range.
\item Result: an economically consistent and \emph{diverse} batch on which the Euler loss is meaningful.
\end{itemize}
\medskip
Compared to DeepHAM (one economy, $N$ agents, $T_v$-step rollout), DEQN runs $N$ economies, each with its own $\mu_n$, evaluating the Euler equation directly each step.
\end{frame}

\begin{frame}{DEQN: distribution transport}
At each epoch, $\mu_{n,t}$ is a histogram over $(a_j, e_k)$. Steps to advance one period:
\medskip

\begin{enumerate}
\item Compute savings $a'_{j,k}$ for every grid point under the current policy.
\item \textbf{Linear interpolation transport}: split mass $\mu(a_j, e_k)$ between the two adjacent grid points with weights $w_{\ell+1}, w_\ell$.
\item \textbf{Markov transition for employment}: distribute mass to $(a_\ell, e')$ weighted by the transition matrix $\Pi_{e\to e'}$.
\item Update $K_{n,t+1} = \sum_{j,k} a_j \cdot \mu_{n,t+1}(a_j, e_k)$.
\end{enumerate}
\medskip

The transport step runs inside \texttt{torch.no\_grad()} so that gradients only propagate through the Euler residual loss, not through the simulation step itself.
\end{frame}

\begin{frame}{DEQN: algorithm outline}
\textbf{Initialisation.} Solve the steady state via SSJ to get $K_{ss}$, asset/employment grids, $\mu_{ss}$, and the Markov transition $\Pi$.
\medskip

\textbf{For each training epoch $t \in \{1,\ldots,T\}$:}
\begin{enumerate}
\item Compute $K_n = \int a\,d\mu_n$ and read $Z_n$ for each cloud copy $n$.
\item Sample $M$ random $(a, e)$ per copy.
\item Forward pass: evaluate $\sigma_\theta \to (a', c)$.
\item Compute Euler residuals $\varepsilon^2$ via expectations over $(Z', e')$.
\item Backward pass: Adam update on $\theta$.
\item Transport: $\mu_n \to \mu_n'$ inside \texttt{no\_grad()}; draw next $Z_n$.
\end{enumerate}
\medskip

\textbf{Convergence.} Loss decreases over $\sim 2000$ epochs; Euler errors below $10^{-3}$ indicate the policy solves the model to high accuracy.
\end{frame}

% Validation frames ported from MiniCourse M2_T4_KS_DEQN section 5.
\begin{frame}{Validation: Euler residuals over training}
\begin{itemize}
\item Track $\log_{10}\!|EE|$ across iterations on a held-out validation sample of $(a, e, Z, K)$.
\medskip
\item Standard targets (Judd 1998, applied to KS):
  \begin{itemize}
  \item Acceptable: $\log_{10}|EE| < -3$
  \item Good: $\log_{10}|EE| < -4$
  \item In this lecture's KS run: \textbf{below $10^{-3}$ over $\sim 2000$ epochs}.
  \end{itemize}
\medskip
\item \textbf{What healthy training looks like}:
  \begin{enumerate}
  \item Loss curve monotone-decreasing on log scale (no oscillation).
  \item Aggregate $K_n$ stays near steady-state range across cloud copies.
  \item Validation Euler error tracks training error (no overfitting --- expected, since the loss \emph{is} the optimality condition).
  \end{enumerate}
\end{itemize}
\end{frame}

\begin{frame}{Validation: economic plausibility}
\begin{itemize}
\item Beyond Euler errors, check that the \emph{economics} is sensible:
  \begin{enumerate}
  \item Savings policy $a'(a, e, Z, K)$ is monotone increasing in $a$ for each fixed $(e, Z, K)$.
  \item Employed agents save more than unemployed at the same $a$ (precautionary motive).
  \item Aggregate $K_n$ tracks the KS log-linear PLM slope-and-intercept reasonably.
  \item Wealth distribution under simulation matches Aiyagari-style steady-state shape (right-skewed, fat upper tail).
  \end{enumerate}
\medskip
\item \textbf{Diagnostic plots} for the lab:
  \begin{itemize}
  \item Loss curve on a log scale.
  \item Policy heatmap $a' = \sigma_\theta(a, e, Z, K)$ at fixed $(Z, K)$.
  \item Stationary wealth histogram.
  \item Cloud trajectory of $K_n$ over training (fan plot).
  \end{itemize}
\end{itemize}
\end{frame}

\begin{frame}{DEQN vs.\ DeepHAM: comparison}
\begin{center}
\small
\begin{tabular}{lll}
\toprule
 & \textbf{DeepHAM} & \textbf{DEQN}\\
\midrule
Networks & Value $+$ policy (2) & Policy only (1)\\
Optimality & Bellman / fictitious play & Euler residual\\
Distribution rep.\ & Generalised moments & Full histogram\\
Training data & $N$ agents, 1 economy & $N$ economies, cloud sim\\
Forward sim & Yes, $T_v$ steps & Yes, 1 step / epoch\\
Distribution evolution & Implicit via moments & Linear-interp transport\\
Scalability & Moments fixed; NN width grows & Histogram sized to grid\\
Differentiability & Through truncated rollout & Euler loss only\\
\bottomrule
\end{tabular}
\end{center}
\medskip
\begin{itemize}
\item Both avoid the curse of dimensionality of the underlying KS problem.
\item DeepHAM has richer distributional info via learned moments; DEQN has a simpler architecture and direct Euler-based training.
\end{itemize}
\end{frame}

\begin{frame}{Recap: the modern HA toolbox}
\begin{itemize}
\item Aiyagari (1994): a non-trivial nested fixed point, classically solved by VFI + bisection on $r$.
\medskip
\item Krusell--Smith (1998): adds aggregate uncertainty, classically tamed by bounded-rationality moments.
\medskip
\item DeepHAM: replaces moments with \emph{learned} sufficient statistics; alternates supervised value learning and fictitious-play policy improvement.
\medskip
\item DEQN: drops the value function entirely; trains a single policy network on Euler residuals over a cloud of simulated economies.
\medskip
\item Both deep-learning methods generalise straightforwardly to:
  \begin{itemize}
  \item Multiple assets (HANK)
  \item OLG with aggregate shocks (next part)
  \item Mean-Field Games in continuous time (Part X)
  \end{itemize}
\end{itemize}
\end{frame}

%=============================================================================
\section{OLG with Aggregate Uncertainty}
% FullCourse-only

\begin{frame}{From infinite-horizon to overlapping generations}
\begin{itemize}
\item In Aiyagari/KS, agents are infinitely lived. Many quantitative questions require finite life cycles:
  \begin{itemize}
  \item Social security and pension reform
  \item Retirement savings and 401(k)
  \item Lifecycle wealth accumulation
  \item Education financing and student debt
  \end{itemize}
\medskip
\item \textbf{OLG model} (Diamond, Auerbach--Kotlikoff): agents live $J$ periods (or cohorts), are born, age, and die.
\medskip
\item State vector now includes the \emph{age} of each cohort, plus its asset/income distribution.
\medskip
\item Standard OLG with aggregate shocks remains computationally hard --- the per-cohort distribution is itself an infinite-dimensional state.
\end{itemize}
\end{frame}

\begin{frame}{OLG with aggregate uncertainty: state vector}
With $J$ cohorts indexed by age $j = 1, \ldots, J$:
\medskip
\begin{itemize}
\item Cohort $j$ has its own asset distribution $\Gamma_j$ over $(a_j, y_j)$.
\medskip
\item Aggregate state: $S = (Z, \Gamma_1, \Gamma_2, \ldots, \Gamma_J)$.
\medskip
\item Aggregate capital and labour:
\[
K_t = \sum_{j=1}^{J} \int a\,d\Gamma_{j,t}, \qquad
L_t = \sum_{j=1}^{J} \int \bar\ell_j\, y\, d\Gamma_{j,t},
\]
where $\bar\ell_j$ is the cohort-specific labour endowment (often hump-shaped).
\medskip
\item Lifecycle Bellman: $V_j(a_j, y_j, S) = \max\,\{u(c) + \beta\, \mathbb{E}[V_{j+1}(a',y',S')]\}$, with $V_{J+1}\equiv 0$ (or a terminal bequest function).
\end{itemize}
\end{frame}

\begin{frame}{OLG state vector: budget constraints and lifecycle profile}
Within each cohort, the household's intertemporal problem at age $j$ has the budget constraint
\[
c_j + a_{j+1} \;=\; (1+r_t)\,a_j \;+\; w_t\,\bar\ell_j\, y_j, \qquad j = 1, \ldots, J,
\]
with terminal condition $a_{J+1} = 0$ (no bequest) or $a_{J+1} = b$ (warm-glow bequest).
\medskip
\begin{itemize}
\item \textbf{Age-profile labour endowment} $\bar\ell_j$ is hump-shaped: rising 25--50, declining 50--65, zero after 65 (retirement).
\item \textbf{Cohort dimensionality.} With $J = 60$ cohorts and a 2-dim individual state $(a, y)$, the full distribution $\{\Gamma_j\}_{j=1}^J$ lives in a $J \times 2 = 120$-dim space.
\item \textbf{Backward induction (deterministic case).} Without aggregate shocks, solve $V_J \to V_{J-1} \to \cdots \to V_1$ cohort-by-cohort; aggregate shocks couple all cohorts simultaneously.
\end{itemize}
\end{frame}

\begin{frame}{Why DL helps for OLG with aggregate shocks}
\begin{itemize}
\item Without aggregate shocks: standard backward induction over cohorts works.
\medskip
\item \textbf{With} aggregate shocks: each cohort's continuation value depends on the future joint distribution of all cohorts $\Rightarrow$ exponential state-space blow-up.
\medskip
\item Deep neural networks let us:
  \begin{itemize}
  \item Approximate $V_j(a, y, S)$ with a single MLP per cohort (or one MLP with cohort as an input).
  \item Encode $S$ via learned summary statistics (analogous to DeepHAM moments).
  \item Train across cohorts in parallel on a GPU.
  \end{itemize}
\medskip
\item Recent reference: Brumm, Kubler \& Scheidegger (2023, working paper), ``Economics-inspired neural networks with stabilizing homotopies''.
\end{itemize}
\end{frame}

\begin{frame}{Homotopy-stabilised training (Brumm--Kubler--Scheidegger 2023)}
\begin{itemize}
\item OLG with aggregate uncertainty is hard to train: borrowing constraints, occasionally-binding bond limits, and discontinuities in the policy create unstable gradients.
\medskip
\item \textbf{Stabilising homotopy}: introduce constraints \emph{gradually} during training.
  \begin{itemize}
  \item Start with a relaxed (smooth) version of the model: no binding constraint, smooth penalties.
  \item Slowly tighten the constraints toward their true value as training progresses.
  \item By the time the constraints are at their true level, the network is already in a good basin.
  \end{itemize}
\medskip
\item Empirically, homotopy is essential for getting OLG-with-aggregate-shocks training to converge.
\medskip
\item Implementation often combines JAX (for automatic differentiation through long simulation rollouts) with the Adam optimiser and a multistep homotopy schedule.
\end{itemize}
\end{frame}

\begin{frame}{Homotopy schedule: results and convergence diagnostics}
A typical 3-stage homotopy for OLG with aggregate uncertainty:
\medskip
\begin{enumerate}
\item \textbf{Stage 1 --- no aggregate shocks.} Fix $Z = \bar Z$; train on the deterministic OLG. Loss converges to $\sim 10^{-5}$ in a few thousand epochs.
\item \textbf{Stage 2 --- small shocks.} Re-introduce aggregate shocks at $\sigma_Z^{(2)} = 0.3\,\sigma_Z^{\text{full}}$. Loss rises to $\sim 10^{-3}$, then converges.
\item \textbf{Stage 3 --- full shocks.} Restore $\sigma_Z^{\text{full}}$; loss stabilises at $\sim 10^{-3}$.
\end{enumerate}
\medskip
\textbf{Convergence diagnostics:}
\begin{itemize}
\item Euler residuals below $10^{-3}$ on a held-out validation set.
\item Cohort-by-cohort policy plots match the Stage-1 deterministic baseline at the modal aggregate state.
\item Cross-stage warm-starting: each stage's final parameters initialise the next.
\end{itemize}
\textit{Why homotopy works:} the policy never sees a discontinuous jump in the loss landscape.
\end{frame}

\begin{frame}{Lab12\_OLG.ipynb preview}
\begin{itemize}
\item \texttt{Notebooks/Lab12\_OLG.ipynb} is a JAX-based implementation of the Brumm--Kubler--Scheidegger OLG model with stabilising homotopies.
\medskip
\item Highlights:
  \begin{itemize}
  \item 7 cohorts spanning a 72-period life cycle, life-cycle labour endowment.
  \item Aggregate TFP shocks following an AR(1).
  \item Capital adjustment costs.
  \item State-dependent depreciation (TFP-tied).
  \item Two-stage homotopy: capital-only model first, then add bond.
  \end{itemize}
\medskip
\item Use this notebook as an \textbf{optional follow-up} after Lab06: students who finish the KS lab early can explore OLG.
\medskip
\item In production research, the same toolkit handles multi-asset HANK and asset-pricing models.
\end{itemize}
\medskip
\emph{Next (T5a):} the continuous-time view --- Mean-Field Games, HJB, and the continuum-of-agents limit.
\end{frame}

\end{document}
